% GENERATED FILE. Edit canonical lessons, then run npm run generate:books.
% canonical-source-hash: fnv1a64:63e7629f0fc135ad
% canonical-lessons: GE-C456-lichtung, GE-C456-gebuesch, GE-C456-farn, GE-C456-moos, GE-C456-rascheln

\chapter{Im Wald}
\label{ch:ge-im-wald}
\hlchaptermodality{456}

\begin{chapteropening}
\textbf{By the end of this chapter:} \emph{I can describe a wood: the clearing, the bushes, the fern, the moss and rustling.}

Complete the last lesson of chapter 456: I can describe a wood: the clearing, the bushes, the fern, the moss and rustling.
\end{chapteropening}

\section[die Lichtung]{die Lichtung — the clearing (in a wood)}

\label{lesson:GE-C456-lichtung}



\begin{quote}
Before the new one: say the German for to splash, then the German for to flow.
\end{quote}

\subsection*{You'll want to know: die Lichtung}
\textbf{die Lichtung} — \textquotedblleft{}the clearing (in a wood)\textquotedblright{}.

\begin{quote}
Auf der Lichtung machen wir ein Picknick.
\end{quote}

\begin{quote}
\emph{We have a picnic in the clearing.}
\end{quote}

Plural: \emph{die Lichtungen}.

\begin{cousinweb}
\emph{Lichtung} is built on \emph{lichten} 'to thin out, to make clear', from \emph{licht} 'light, open', inherited from Proto-Germanic \emph{leuhtą} 'light' (EtymDB).
\end{cousinweb}

\begin{grammarlens}[title={What you've built}]
One more word for telling, explaining and describing at B1.
\end{grammarlens}

\subsection*{Your turn}
\begin{itemize}
\raggedright
  \item \emph{Say it:} \emph{die Lichtung}
  \item \emph{Say it:} \emph{die Lichtung}, once more
  \item \emph{Say it:} a sentence of your own with \emph{die Lichtung}
  \item \emph{Recall:} say the German for to splash, then the German for to flow, then say \emph{die Lichtung} again
\end{itemize}

\begin{tcolorbox}[breakable,colback=teal!4,colframe=teal!35!black,title={Before you move on}]
What does die Lichtung mean? (\textquotedblleft{}The clearing (in a wood)\textquotedblright{}.) Say it once more.
\end{tcolorbox}
\section[das Gebüsch]{das Gebüsch — the bushes, the undergrowth}

\label{lesson:GE-C456-gebuesch}



\begin{quote}
Before the new one: say the German for to flow, then the German for the clearing.
\end{quote}

\subsection*{You'll want to know: das Gebüsch}
\textbf{das Gebüsch} — \textquotedblleft{}the bushes, the undergrowth\textquotedblright{}.

\begin{quote}
Ein Kaninchen versteckt sich im Gebüsch.
\end{quote}

\begin{quote}
\emph{A rabbit hides in the bushes.}
\end{quote}

Plural: \emph{die Gebüsche}. One bush is \emph{der Busch}.

\begin{cousinweb}
\emph{Gebüsch} is built from \emph{ge-}, which groups things together, and \emph{Busch} 'bush', inherited from Old High German \emph{busk} (EtymDB); English \emph{bush} is a cousin.
\end{cousinweb}

\begin{grammarlens}[title={What you've built}]
One more word for telling, explaining and describing at B1.
\end{grammarlens}

\subsection*{Your turn}
\begin{itemize}
\raggedright
  \item \emph{Say it:} \emph{das Gebüsch}
  \item \emph{Say it:} \emph{das Gebüsch}, once more
  \item \emph{Say it:} a sentence of your own with \emph{das Gebüsch}
  \item \emph{Recall:} say the German for to flow, then the German for the clearing, then say \emph{das Gebüsch} again
\end{itemize}

\begin{tcolorbox}[breakable,colback=teal!4,colframe=teal!35!black,title={Before you move on}]
What does das Gebüsch mean? (\textquotedblleft{}The bushes, the undergrowth\textquotedblright{}.) Say it once more.
\end{tcolorbox}
\section[der Farn]{der Farn — the fern}

\label{lesson:GE-C456-farn}



\begin{quote}
Before the new one: say the German for the clearing, then the German for the bushes.
\end{quote}

\subsection*{You'll want to know: der Farn}
\textbf{der Farn} — \textquotedblleft{}the fern\textquotedblright{}.

\begin{quote}
Unter den Bäumen wächst grüner Farn.
\end{quote}

\begin{quote}
\emph{Green ferns grow under the trees.}
\end{quote}

Plural: \emph{die Farne}.

\begin{cousinweb}
\emph{Farn} is an inherited Germanic word; EtymDB lists English \emph{fern} and Dutch \emph{varen} as cousins and links them to an old word for a feather or wing.
\end{cousinweb}

\begin{grammarlens}[title={What you've built}]
One more word for telling, explaining and describing at B1.
\end{grammarlens}

\subsection*{Your turn}
\begin{itemize}
\raggedright
  \item \emph{Say it:} \emph{der Farn}
  \item \emph{Say it:} \emph{der Farn}, once more
  \item \emph{Say it:} a sentence of your own with \emph{der Farn}
  \item \emph{Recall:} say the German for the clearing, then the German for the bushes, then say \emph{der Farn} again
\end{itemize}

\begin{tcolorbox}[breakable,colback=teal!4,colframe=teal!35!black,title={Before you move on}]
What does der Farn mean? (\textquotedblleft{}The fern\textquotedblright{}.) Say it once more.
\end{tcolorbox}
\section[das Moos]{das Moos — the moss}

\label{lesson:GE-C456-moos}



\begin{quote}
Before the new one: say the German for the bushes, then the German for the fern.
\end{quote}

\subsection*{You'll want to know: das Moos}
\textbf{das Moos} — \textquotedblleft{}the moss\textquotedblright{}.

\begin{quote}
Die alten Steine sind voll mit Moos.
\end{quote}

\begin{quote}
\emph{The old stones are covered in moss.}
\end{quote}

Plural: \emph{die Moose}. In slang, \emph{Moos} also means money.

\begin{cousinweb}
\emph{Moos} is inherited through Middle High German \emph{mos} from Old High German \emph{mos} (EtymDB), which also meant a bog; Old English \emph{mos} 'bog, moss' is a cousin.
\end{cousinweb}

\begin{grammarlens}[title={What you've built}]
One more word for telling, explaining and describing at B1.
\end{grammarlens}

\subsection*{Your turn}
\begin{itemize}
\raggedright
  \item \emph{Say it:} \emph{das Moos}
  \item \emph{Say it:} \emph{das Moos}, once more
  \item \emph{Say it:} a sentence of your own with \emph{das Moos}
  \item \emph{Recall:} say the German for the bushes, then the German for the fern, then say \emph{das Moos} again
\end{itemize}

\begin{tcolorbox}[breakable,colback=teal!4,colframe=teal!35!black,title={Before you move on}]
What does das Moos mean? (\textquotedblleft{}The moss\textquotedblright{}.) Say it once more.
\end{tcolorbox}
\section[rascheln]{rascheln — to rustle}

\label{lesson:GE-C456-rascheln}



\begin{quote}
Before the new one: say the German for the fern, then the German for the moss.
\end{quote}

\subsection*{You'll want to know: rascheln}
\textbf{rascheln} — \textquotedblleft{}to rustle\textquotedblright{}.

\begin{quote}
Im Herbst raschelt das Laub unter unseren Füßen.
\end{quote}

\begin{quote}
\emph{In autumn the leaves rustle under our feet.}
\end{quote}

\emph{Es raschelt im Gebüsch} — something is rustling in the bushes.

\begin{cousinweb}
\emph{rascheln} is a sound word that imitates the noise of dry leaves or paper (EtymDB lists the sense 'to rustle').
\end{cousinweb}

\begin{grammarlens}[title={What you've built}]
One more word for telling, explaining and describing at B1.
\end{grammarlens}

\subsection*{Your turn}
\begin{itemize}
\raggedright
  \item \emph{Say it:} \emph{rascheln}
  \item \emph{Say it:} \emph{rascheln}, once more
  \item \emph{Say it:} a sentence of your own with \emph{rascheln}
  \item \emph{Recall:} say the German for the fern, then the German for the moss, then say \emph{rascheln} again
\end{itemize}

\begin{tcolorbox}[breakable,colback=teal!4,colframe=teal!35!black,title={Before you move on}]
What does rascheln mean? (\textquotedblleft{}To rustle\textquotedblright{}.) Say it once more.
\end{tcolorbox}
